\chapter{روش پیشنهادی}
\label{chp:proposed_method}
\section{مقدمه}
در این فصل سعی داریم با ارائه توضیحاتی جامع به تشریح روش پیشنهادی برای انجام آزمایش‌ها بپردازیم. در ابتدا با توضیحات مربوط به مجموعه داده‌های استفاده شده در بخش \ref{sec:data} با مجموعه داده‌های استفاده شده آشنا می‌شویم. سپس در بخش \ref{sec:labeling} به شیوه برچسب‌گذاری توییت‌ها و فرایند بهینه‌سازی متغییر‌های روش سه‌مانع می‌پردازیم. در ادامه و بخش \ref{sec:fine-tuning} به شیوه انجام تنظیم دقیق و مدل‌های هدف در پژوهش می‌رسیم و با جزيیات این فرایند آشنا می‌شویم. در بخش \ref{sec:prompt-tuning} به شیوه تزریق اطلاعات بازار و اطلاعات زمانی به توییت ها می‌پردازیم. سپس در بخش \ref{sec:signal-generation} به روش‌های محاسبه حجم معاملات و سیگنال روزانه می‌پردازیم. و در انتها با بخش \ref{sec:backtest} به ارائه روش‌های پیشنهادی برای استراتژی‌های مناسب به مدل‌های زبانی و معیار‌های استفاده شده اشاره می‌کنیم.  
\section{جمع‌آوری داده}
برای پشتیبانی از پیش‌بینی روند بازار با استفاده از احساسات و شاخص‌های فنی، مجموعه داده‌هایی شامل توییت‌های روزانه، شاخص‌های فنی و رویدادهای مهم بیتکوین جمع‌آوری کردیم. هر مجموعه داده هدف مشخصی دارد: توییت‌ها برای آموزش و ارزیابی مدل استفاده شدند؛ شاخص‌های فنی به عنوان ویژگی‌های مبتنی بر قیمت در مدل‌های طبقه‌بندی و برای افزودن زمینه بازار در فرآیند تنظیم پرسش به کار رفتند؛ و رویدادهای تاریخی برای ارزیابی عملکرد مدل در تاریخ‌های کلیدی، به‌ویژه در دوره‌های نوسان شدید بازار یا چرخه‌های خبری تأثیرگذار، مورد استفاده قرار گرفتند.
\label{sec:data}
\begin{itemize}
\item \textbf{مجموعه توییت‌ها:}
در این پژوهش از داده‌های متنی توییتر برای تحلیل احساسات و پیش‌بینی حرکات کوتاه‌مدت قیمت بیتکوین استفاده شد. توییت‌ها شامل دو مجموعه داده متفاوت است که شامل توییت‌های جمع‌آوری شده از سال ۲۰۱۵ تا فوریه ۲۰۲۱ بوده \cite{prebit_zou_2023} و دومین مجموعه داده توییت‌های بیتکوین شامل داده‌هایی از شروع سال ۲۰۲۱ تا ژانویه سال ۲۰۲۳ است که توییت‌های روزانه مربوط به بیتکوین را در بر می‌گیرد. برای اطمینان از کیفیت داده‌ها، چندین مرحله پیش‌پردازش \cite{hickman2022text} بر روی داده‌های متنی انجام شد.

\item \textbf{رویداد‌های تاریخی بیتکوین:}
علاوه بر داده‌های قیمتی و توییت‌ها، از مجموعه داده‌ای شامل ۲۲۷ رویداد تاریخی از سال ۲۰۰۹ تا ۲۰۲۴ که تأثیر قابل‌توجهی بر بازار بیتکوین داشتند استفاده شد. این مجموعه داده برای نمونه‌برداری رویداد\LTRfootnote{Event Sampling} از روزهای خاصی که اتفاقات تاثیر‌گذار در روند قیمت بیتکوین رخ داده بود، به منظور ارزیابی مدل در مواجهه با این نوع روزها استفاده شد \cite{deep_ding_2015}. این مجموعه داده فارغ از برچسب روز یا محل اتفاق در پنجره برچسب‌ها تهیه شده و از این حیث به ارزیابی بدون تبعیض کمک می‌کند. این مجموعه داده در \href{https://www.kaggle.com/datasets/hamidmoradi21/bitcoin-historical-events}{دسترس عموم} قرار گرفته است.
\item \textbf{داده‌های قیمتی و شاخص‌های فنی}
برای تحلیل حرکات قیمتی بیتکوین، از داده‌های تاریخی قیمت بیتکوین استفاده شد. این داده‌ها شامل قیمت‌های پایانی روزانه از سال ۲۰۱۵ تا ۲۰۲۳ بود. به منظور استخراج ویژگی‌های لازم برای مدل‌سازی، قیمت‌ها در کنار شاخص‌های فنی مرتبط مانند شاخص قدرت نسبی و نرخ تغییرات مورد استفاده قرار گرفتند.

این دو شاخص فنی اصلی برای تحلیل وضعیت بازار بیتکوین محاسبه شدند

\begin{itemize}
    \item شاخص قدرت نسبی: که بر اساس تغییرات قیمت در دوره‌های زمانی مختلف محاسبه شده و فرمول مربوطه \ref{eq:rsi} برای پیاده‌سازی آن استفاده شده. محدوده‌های ۳۰ و ۷۰ به عنوان مرزهای اشباع فروش و اشباع خرید در نظر گرفته شده‌اند که به عنوان مرز طبقه‌بندی صعودی و نزولی استفاده شده.

\begin{equation} \label{eq:rsi} \text{\lr{RSI}} = 100 - \left( \frac{100}{1 + \frac{\text{\lr{Average Gain}}}{\text{\lr{Average Loss}}}} \right) \end{equation}

    \item \textbf{نرخ تغییرات:} این شاخص یکی از ابزارهای مهم تحلیل تکنیکال است که تغییرات نسبی قیمت در یک بازه زمانی مشخص را نشان می‌دهد. نرخ تغییرات با استفاده از فرمول \ref{eq:ROC} محاسبه می‌شود که در آن \( P_{\text{close}}(d_i) \) قیمت پایانی در روز \( d_i \) و \( P_{\text{close}}(d_{i-1}) \) قیمت پایانی در روز قبل (\( d_{i-1} \)) است. این فرمول به صورت زیر تعریف می‌شود:

\begin{equation} \label{eq:ROC}
ROC_i = \frac{P_{\text{close}}(d_i) - P_{\text{close}}(d_{i-1})}{P_{\text{close}}(d_{i-1})} \times 100 
\end{equation}

فرمول بالا میزان تغییرات روزانه قیمت را به صورت درصدی محاسبه می‌کند و اطلاعات ارزشمندی درباره روند قیمت ارائه می‌دهد. برای تفسیر این شاخص، می‌توان سیگنال‌هایی را تعریف کرد. اگر \( ROC_i > 0 \)، نشان‌دهنده یک روند صعودی است که قیمت افزایش یافته است. اگر \( ROC_i < 0 \)، به معنای روند نزولی و کاهش قیمت است. در نهایت، اگر \( ROC_i \approx 0 \)، تغییرات قابل‌توجهی رخ نداده و سیگنال خنثی در نظر گرفته می‌شود. این شاخص نقش مهمی در تحلیل تغییرات کوتاه‌مدت بازار ایفا می‌کند و می‌تواند به صورت مستقیم برای تولید سیگنال‌های معاملاتی و پیش‌بینی روند بازار مورد استفاده قرار گیرد.

\end{itemize}
\end{itemize}

\subsection{پیش‌پردازش}
برای اطمینان از کیفیت و یکنواختی داده‌ها، چندین مرحله پیش‌پردازش\LTRfootnote{Preprocessing} بر روی داده‌های توییت‌های مالی اعمال می‌شود. این گام‌ها شامل عادی‌سازی\LTRfootnote{Normalization} متن، استخراج ایموجی‌ها و هشتگ‌ها\LTRfootnote{Hashtags}، فیلتر کردن تبلیغات و ایردراپ‌ها\LTRfootnote{Airdrops}، و یکسان‌سازی نام دارایی‌ها است. یکی از اولین گام‌های مهم در پیش‌پردازش توییت‌ها، عادی‌سازی و تمیز کردن متن است. در این مرحله، ابتدا تمامی متون به حروف کوچک تبدیل می‌شوند تا از تفاوت‌های حروف بزرگ و کوچک جلوگیری شود. سپس، نشانی‌های اینترنتی، نشانگر\LTRfootnote{ID}کاربران و علائم نگارشی که در تحلیل احساسات\LTRfootnote{Sentiment Analysis} معمولاً کاربردی ندارند، حذف می‌شوند. این مراحل به کاهش پراکندگی\LTRfootnote{Noise} در داده‌ها و بهبود کیفیت داده‌های ورودی مدل کمک می‌کنند \cite{twitter_pak_2010}. یکی از ویژگی‌های بارز در توییت‌ها، استفاده گسترده از ایموجی‌ها\LTRfootnote{Emojis} و هشتگ‌هاست. ایموجی‌ها غالباً بیانگر احساسات کاربران هستند و می‌توانند به عنوان شاخصی برای تعیین نوع احساسات استفاده شوند. همچنین، هشتگ‌ها می‌توانند محتوای مهمی در مورد موضوع توییت و فضای کلی آن فراهم کنند. استخراج این عناصر و پردازش جداگانه آن‌ها به غنی‌تر شدن مدل‌های تحلیل احساسات کمک می‌کند \cite{emojipowered_chen_2018}

در توییت‌های مربوط به بازارهای مالی، تعداد زیادی تبلیغ\LTRfootnote{Advertisements} و رویدادهایی نظیر ایردراپ‌ها وجود دارند که این توییت‌ها عموماً ارزش تحلیلی کمی دارند. فیلتر کردن کلمات کلیدی مرتبط با تبلیغات و ایردراپ‌ها باعث حذف این توییت‌ها از مجموعه داده شده و در نهایت به بهبود کیفیت داده‌ها منجر می‌شود. در داده‌های توییتر مرتبط با بازارهای مالی، نام دارایی‌ها به اشکال مختلفی نوشته می‌شوند. برای مثال، بیت‌کوین ممکن است به صورت‌های مختلفی نظیر (\textit{\lr{‌Bitcoin}}, \textit{\lr{BTC}}, و یا \textit{بیتکوین}) نوشته شود. یکی از گام‌های مهم در پیش‌پردازش داده‌های مالی، یکسان‌سازی این نام‌ها و جایگزینی آن‌ها با یک نماد واحد است. این کار از پراکندگی داده‌ها جلوگیری کرده و دقت مدل‌های تحلیل احساسات را بهبود می‌بخشد. به منظور بهبود توانایی مدل در شناخت کلمات مرتبط و افزایش تعمیم‌پذیری آن، فرآیند ریشه‌یابی\LTRfootnote{Lemmatization} روی متون اعمال می‌شود. ریشه‌یابی به مدل اجازه می‌دهد تا شکل‌های مختلف یک کلمه مانند \textit{\lr{Running}} و  \textit{\lr{Ran}} را به عنوان یک مفهوم واحد در نظر بگیرد. این امر حجم ویژگی‌های ورودی مدل را کاهش داده و دقت آن را در تشخیص مفاهیم مرتبط افزایش می‌دهد. در ادامه به بعضی از جزئیات پیاده سازی این روش می‌پردازیم.


\section{برچسب‌گذاری بر پایه قیمت}
\label{sec:labeling}
هدف این بخش از تحقیق، پیش‌بینی حرکات کوتاه‌مدت بازار از طریق طبقه‌بندی توییت‌ها بر اساس تأثیر کوتاه‌مدت آن‌ها با استفاده از روش برچسب‌گذاری سه‌مانع است. این روش بر روی تعیین مرزهای زمانی و سود/زیان برای هر توییت تمرکز دارد و سعی دارد معنای دقیق و قابل بازتولید برای احساسات نزولی، صعودی و خنثی ایجاد کند. فرضیه‌ی اصلی این است که استفاده از برچسب‌گذاری بازار محور، می‌تواند پیش‌بینی‌های بهتری برای روند ‌کوتاه‌ مدت نسبت به تحلیل سنتی احساسات ارائه دهد، که غالباً تعاریف مبهمی از احساسات نزولی و صعودی دارد. ما از روش برچسب‌گذاری سه‌مانع \cite{advances_deprado_2018} برای برچسب‌گذاری داده‌های قیمت روزانه استفاده می‌کنیم که به عنوان برچسب کوتاه‌ مدت توییت‌ها استفاده می‌شود. شیوه کلی استفاده از این روش در بخش \ref{sec:tbl} آمده و تشریح شده است. با توجه به تغییرات اساسی بازار در طول زمان انتظار ما از حرکات بازار نیز تغییر می‌کند \cite{ang2012regime}. از این رو استفاده از مرز‌های ثابت و از پیش تعیین شده برای برچسب‌گذاری داده‌های بازار مناسب نیست. در این بخش می‌خواهیم به شیوه محاسبه حدود به صورت پویا با استفاده از نوسانات تاریخی بازار و بهینه‌سازی متغییر‌های این روش با استفاده از نسبت شارپ به دست آمده از پیش‌آزمون بپردازیم \cite{zhang2020deep}.تأثیر احساسات موجود در رسانه‌های اجتماعی بر بازارهای مالی اغلب تا ۱۵ روز ادامه دارد، اما بسیاری از مطالعات از آستانه‌های یک‌روزه برای برچسب‌گذاری روندها استفاده می‌کنند \cite{twitter_vallecruz_2021}. این ساده‌سازی، اثرات احساسات در روند کوتاه‌ مدت را در نظر نمی‌گیرد و با شیوه‌های واقعی معاملات که روندها را در چندین روز دنبال می‌کنند، سازگار نیست. برچسب‌گذاری مؤثر باید این پویایی‌ها را در نظر بگیرد و همچنین روش‌هایی مانند حد سود، حد ضرر و حد زمانی را برای مدیریت ریسک و هم‌راستا کردن پیشبینی‌ها با استراتژی‌های عملیاتی، یکپارچه کند. برای پاسخ به این چالش‌ها، روش برچسب‌گذاری مبتنی بر بازار پیشنهاد شده است که از چارچوب سه‌مانع الهام گرفته است. این روش از نوسانات تاریخی و حدود بازارمحور برای ایجاد برچسب‌های توضیح‌پذیر و قابل اجرا برای مدل‌های یادگیری ماشین و عوامل معاملاتی بهره می‌برد.

\begin{enumerate}
    \item \textbf{برآورد نوسان تاریخی:}
    نوسان تاریخی\LTRfootnote{Historical Volatility} معیاری است که میزان تغییرات قیمت یک دارایی را در طول یک دوره زمانی مشخص اندازه‌گیری می‌کند. این نوسان با استفاده از داده‌های قیمت گذشته محاسبه می‌شود و به عنوان یک شاخص از ریسک و عدم قطعیت در بازار مورد استفاده قرار می‌گیرد. برای محاسبه نوسان تاریخی، معمولاً از انحراف معیار بازده‌های لگاریتمی\LTRfootnote{Log Returns} استفاده می‌شود که تغییرات نسبی قیمت را نشان می‌دهند. در برچسب‌گذاری پیشنهاد شده، نوسان تاریخی به عنوان مبنایی برای تنظیم حدود صعودی و نزولی به صورت پویا استفاده می‌شود. به این صورت که حدود به عنوان ضریبی از نوسان در ابتدای هر بازه زمانی تنظیم می‌شوند، که این امر باعث می‌شود تا حدود به تغییرات بازار واکنش نشان دهند و برچسب‌ها به صورت دقیق‌تری اختصاص داده شوند. این رویکرد به ما اجازه می‌دهد تا تأثیرات متنی را به تغییرات واقعی بازار مرتبط کنیم و تحلیل‌های دقیق‌تری ارائه دهیم 
    نوسان تاریخی بازده‌های لگاریتمی با استفاده از روش میانگین متحرک نمایی وزنی\LTRfootnote{Exponentially Weighted Moving Average} تخمین زده می‌شود. بازده‌های لگاریتمی به صورت زیر محاسبه می‌شوند:
    \begin{align}
    r_t &= \ln\left(\frac{P_t}{P_{t-1}}\right) \\
    \sigma_t &= \sqrt{\frac{\alpha}{1 - (1 - \alpha)^\tau} \sum_{i=1}^\infty (1-\alpha)^{i-1} (r_{t-i})^2}
    \end{align}
    که در آن \(\alpha = \frac{2}{\tau + 1}\) عامل هموارسازی و \(\tau\) طول پنجره است. نوسان میانگین متحرک نمایی وزنی به عنوان انحراف معیار بازده‌های لگاریتمی محاسبه می‌شود که تغییرات قیمت اخیر را سنگین‌تر در نظر می‌گیرد تا شرایط پویای بازار را بهتر شبیه‌سازی کند. در این پژوهش، طول پنجره \(\tau\) برابر با ۳۰ روز تنظیم شده است.

    \item \textbf{محاسبه حدود و تخصیص برچسب:}  
    حدود بالا و پایین ابتدا با استفاده از نوسان تاریخی محاسبه می‌شوند تا حرکات قیمت برچسب‌گذاری شوند، به گونه‌ای که پیش‌بینی‌ها بدون سوگیری پیش‌بینی آینده باشند. حدود برای زمان \(i\) به صورت زیر تعریف می‌شوند:
    \begin{align}
    \text{U}_i &= P_i + (P_i \times \sigma_i \times F_\text{u}) \label{eq:upper} \\
    \text{L}_i &= P_i - (P_i \times \sigma_i \times F_\text{l}) \label{eq:lower}
    \end{align}
    که در آن \(P_i\) قیمت بسته‌شده، \(\sigma_i\) نوسان تاریخی و \(F_\text{u}\)، \(F_\text{l}\) ضرایب حدود بالا و پایین هستند که بر اساس پویایی‌های بازار بهینه‌سازی می‌شوند. مانع عمودی (\(\text{V}_i\)) به عنوان حداکثر افق زمانی برای برچسب‌گذاری تعریف می‌شود که بین ۸ تا ۱۵ روز تغییر می‌کند (این مقادیر به طور معمول بازه تاثیر‌گذاری توییت‌هاست ولی در خصوص تمام توییت‌ها صدق نمی‌کند \cite{twitter_vallecruz_2021}). تخصیص برچسب بر اساس قواعد زیر انجام می‌شود:
    \begin{itemize}
        \item برچسب \textit{صعودی} زمانی تخصیص داده می‌شود که \(P_t \geq \text{U}_i\) قبل از رسیدن به \(\text{V}_i\).
        \item برچسب \textit{نزولی} زمانی تخصیص داده می‌شود که \(P_t \leq \text{L}_i\) قبل از رسیدن به \(\text{V}_i\).
        \item برچسب \textit{خنثی} زمانی تخصیص داده می‌شود که نه \(\text{U}_i\) و نه \(\text{L}_i\) تا زمان رسیدن به \(\text{V}_i\) عبور نکنند.
    \end{itemize}

    \item \textbf{بهینه‌سازی و استراتژی معاملاتی:}  
    این روش سه متغیر ضرایب حدود بالا و پایین (\(F_\text{u}, F_\text{l}\)) و طول مانع عمودی (\(\text{V}_i\))را طی دوره‌های شش‌ماهه با استفاده از نسبت شارپ بهینه‌سازی می‌کند. نسبت شارپ به صورت زیر محاسبه می‌شود:
    \begin{align}
    \text{\lr{Sharpe Ratio}} = \frac{\mathbb{E}[R_p - R_f]}{\sigma_p}
    \label{eq:sharpe}
    \end{align}
    که در آن \(R_p\) بازده پرتفوی، \(R_f\) نرخ بدون ریسک (برابر با ۲٪ سالانه) و \(\sigma_p\) انحراف معیار بازده‌ها است. سال ۳۶۵ روزه فرض می‌شود. بهینه‌سازی تضمین می‌کند که حدود (\(\text{U}_i\)، \(\text{L}_i\)، و \(\text{V}_i\)) فقط در صورتی معتبر باشند که قیمت بالا یا پایین تا زمان \(\text{V}_i\) لمس شوند. حداقل مدت زمان روند دو روزه نیز برای افزایش مقاومت برچسب‌گذاری اعمال می‌شود.
\end{enumerate}

این فرآیند برچسب‌گذاری با چارچوب زمانی معمول اثر رسانه‌های اجتماعی بر بازار هم‌راستا است و یک پنجره تصمیم‌گیری مؤثر برای عوامل معاملاتی فراهم می‌کند. این روش به چالش‌های پیش‌بینی در پنجره‌های اولیه پرداخته و با پیشرفت پنجره، قابل‌اعتمادتر می‌شود. شکل \ref{fig:tbl_vis} این روش را نشان می‌دهد که خطوط سبز و قرمز به ترتیب نشان‌دهنده حدود بالا و پایین و برچسب‌های تخصیص‌یافته هستند. در این بخش، لازم است به یک نکته مهم درباره ارتباط بین برچسب‌گذاری بر پایه روند و پیش‌بینی روز بعد اشاره شود.


\section{طبقه‌بندی توییت‌ها بر پایه روند کوتاه‌مدت}
\label{sec:fine-tuning}
در این مطالعه، ما ادعا می‌کنیم که توییت‌ها را بر اساس روند کوتاه‌مدت بازار برچسب‌گذاری می‌کنیم. با این حال، روش برچسب‌گذاری ما، برچسب هر توییت در یک روز معین را به‌عنوان روند روز بعد، که از روش پیشنهادی مبتنی بر روش برچسب‌گذاری پیشنهاد شده به دست آمده است، اختصاص می‌دهد. این موضوع این سوال را مطرح می‌کند: چگونه می‌توانیم ادعای برچسب‌گذاری روند را داشته باشیم در حالی که برچسب‌ها تنها بر اساس روند روز بعد تعیین می‌شوند؟ پاسخ در ماهیت روش مبتنی بر روش برچسب‌گذاری پیشنهادی نهفته است. برچسب روند روز بعد تنها به حرکت بازار در روز بعد وابسته نیست؛ بلکه نشان‌دهنده‌ی پنجره کوتاه‌مدت پیش‌رو (یا روندی) است که پس از آن روز شکل می‌گیرد. دلیل این امر این است که برچسب بر اساس این که کدام یک از موانع سه‌گانه (حد بالا، حد پایین، یا زمان) در بازه مشخص‌شده فعال می‌شود، تعیین می‌گردد. بنابراین، زمانی که برچسب روند روز بعد را به یک توییت اختصاص می‌دهیم، به‌طور ذاتی اطلاعات مربوط به روند کوتاه‌مدت را در نظر می‌گیریم و نه صرفاً تغییر بازار در روز بعد را. در نتیجه در ادامه این پژوهش با اشاره به برچسب روز بعد، در واقع به برچسب روند کوتاه‌مدت اشاره داریم.

\subsection{معرفی مدل‌های طبقه‌بندی}
\label{sec:prompt-tuning}

\begin{figure}[h!]
\centering
\includegraphics[width=1\linewidth]{figures/context-aware-diagram.png}
\caption{فرآیند تنظیم دقیق مدل‌ها}
\label{fig:fine_tuning_process}
\end{figure}

در این مطالعه، سه مدل طبقه‌بندی‌کننده‌ی متن (در این پژوهش صرفا از توییت استفاده شده) معرفی می‌شوند که با هدف پیش‌بینی روندهای کوتاه‌مدت بازار، بر اساس یک روش برچسب‌گذاری نوآورانه طراحی شده‌اند. هدف اصلی این بخش، بررسی قابلیت مدل‌های زبانی\LTRfootnote{Language Model} در پیش‌بینی حرکات بازار به صورت مستقیم از توییت‌ها و مقایسه‌ی این عملکرد با مدل‌های سنتی مبتنی بر تحلیل احساسات است. ابتدا، توانایی مدل‌های زبانی در پیش‌بینی مستقیم روندهای بازار ارزیابی می‌شود تا مشخص شود آیا توییت‌ها دارای ویژگی‌های قابل یادگیری مرتبط با حرکات آتی بازار هستند یا خیر. برای این منظور، مدل‌های سنتی تحلیل احساسات مانند \lr{FinBERT}\cite{prebit_zou_2023} و \lr{CryptoBERT}\cite{sentiment_kulakowski_2023}، با داده‌های دارای برچسب‌های روند بازار بررسی می‌شوند. این گام به عنوان پایه‌ای برای درک اینکه آیا احساسات به تنهایی می‌توانند تغییرات بازار را پیش‌بینی کنند، در نظر گرفته شده است. تنظیم پرسش\LTRfootnote{Prompt-tuning} روشی است که به ما اجازه می‌دهد اطلاعات اضافی مانند اطلاعات بازار و اطلاعات زمانی را به توییت‌ها اضافه کنیم. این روش به مدل کمک می‌کند تا با استفاده از این اطلاعات اضافی، تحلیل دقیق‌تری از توییت‌ها داشته باشد. در ادامه، به بررسی رویکردی می‌پردازیم که اطلاعات بازار را به پرسش‌ها (توییت‌ها) اضافه می‌کند تا تأثیرات بازار را در تحلیل مدل لحاظ کنیم. هدف از این کار، بهبود پیش‌بینی‌های کوتاه‌مدت با اضافه کردن اطلاعاتی نظیر برچسب قبلی (وضعیت بازار در پنجره‌ی زمانی پیشین)، شاخص قدرت نسبی، و نرخ تغییرات به هر توییت به صورت توصیفی است. این اطلاعات با استفاده از الگوریتم‌های محاسباتی بر اساس داده‌های تاریخی قیمت محاسبه می‌شوند. 


\begin{figure}
    \centering
    \includegraphics[width=0.75\linewidth]{figures/tweets.png}
    \caption{توییت \lr{A} قبل از اضافه شدن اطلاعات، توییت \lr{B} بعد از اضافه شدن اطلاعات بازار، و توییت \lr{C} بعد از اضافه شدن اطلاعات زمانی است.}
    \label{fig:tweets}
\end{figure}

هدف از افزودن اطلاعات بازار به توییت‌ها این است که مدل زبانی بتواند یک توییت را نه تنها بر اساس محتوای آن، بلکه بر اساس محیطی که در آن ارسال شده است تحلیل کند. تصویری از شکل کلی استفاده از این روش‌ها در پژوهش فعلی در شکل \ref{fig:fine_tuning_process} آمده است. با توجه به مطالعات فراوانی که نشان می‌دهند افزودن داده‌های قیمتی به تحلیل متنی قدرت پیش‌بینی روند بازار را بهبود می‌بخشد، منطقی است که این اطلاعات به توییت‌ها اضافه شوند\cite{man_2019_financial}. از آنجا که ما به طور مداوم اطلاعات بازار را به توییت‌ها اضافه می‌کنیم، ممکن است مدل به برخی از این اطلاعات بیش‌برازش\LTRfootnote{Overfit} کند و نتواند به راحتی از اطلاعات ارائه شده برای توییت‌های یک روز خاص عبور کند. با توجه به مطالعات مرتبط، افزودن اطلاعات زمانی به مدل زبان به صورت متنی می‌تواند به مدل کمک کند تا اطلاعات کافی در مورد زمینه زمانی توییت دریافت کند. در گام بعدی، مدل \lr{CryptoBERT} با استفاده از برچسب‌های مبتنی بر بازار، برای ساخت مدلی به نام \lr{Context-UnAware} (\lr{CUA}) آموزش داده شد. هدف این مدل، شناسایی الگوهایی در محتوای توییت‌ها است که با روندهای بازار مرتبط هستند، در حالی که دانش موجود مدل در زمینه اصطلاحات رمزارزی و ساختار توییت‌ها حفظ می‌شود. فرآیند آموزش شامل فریز کردن ۱۱ لایه اول \lr{CryptoBERT} برای نگهداشت دانش دامنه‌ای و تغییر تنها سر طبقه‌بندی است. آموزش مدل طی دو اپوک و با استفاده از اعتبارسنجی متقابل ۵‌-گانه انجام شده است تا از بیش‌برازش جلوگیری شود.

در این فرآیند هیچگونه تغییر ساختاری در مدل ایجاد نشده اما برای تنظیم دقیق از بهینه‌ساز \lr{AdamW} با نرخ یادگیری $10^{-5}$ و اندازه دسته\LTRfootnote{‌Batch Size} ۱۲ استفاده شد. همچنین، یک زمان‌بند نرخ یادگیری با فاز گرم شدن خطی (۱۰ درصد مراحل) برای یادگیری تدریجی به کار گرفته شد. بر اساس نتایج تحقیقات پیشین در خصوص هم‌افزایی میان متن و زمینه‌ی بازار \cite{man_2019_financial}، اطلاعات بازار از طریق تنظیم پرسش به ورودی مدل اضافه شد. این تکنیک که در پیش‌بینی مالی و تحلیل احساسات مؤثر بوده است \cite{leveraging_ahmed_2024}، از طریق تنظیم پرسش توان پیش‌بینی مدل‌های زبانی از پیش آموزش‌دیده را بدون تغییر در معماری آن‌ها افزایش می‌دهد. ویژگی‌های اضافه‌شده به هر توییت شامل شاخص قدرت نسبی، نرخ تغییرات، و روند پنجره‌ی زمانی قبلی است. این ویژگی‌ها به دلیل همبستگی بالاتر آن‌ها با روندهای بازار انتخاب شدند (همبستگی $0.26$ برای روند قبلی، $0.14$ برای نرخ بازدهی و $0.13$ برای شاخص قدرت نسبی). این شاخص‌ها به صورت اصطلاحات توصیفی تبدیل شده و پرسش‌هایی به این صورت ایجاد می‌کنند:
\begin{quote}
برچسب قبلی: صعودی، \lr{ROC}: خنثی، \lr{RSI}: نزولی، توییت: بیت‌کوین در حال افزایش است!
\end{quote}

مدلی که بر اساس این توصیفات غنی‌شده ساخته شده، \lr{Context-Aware} (\lr{CA}) نام دارد که نسخه‌ی پیشرفته‌ی مدل \lr{CUA} است. این مدل برای تحلیل دقیق‌تر روندهای بازار با توجه به محتوای توییت‌ها و ویژگی‌های بازار آموزش داده شده است. علاوه بر این، برای درک بهتر اثرات زمانی، مدل با افزودن تاریخ هر توییت به ورودی بهبود داده شد. مدل‌های زبانی معمولاً فاقد آگاهی ذاتی از زمان هستند \cite{timeaware_dhingra_2021, temporal_agarwal_2021}. با افزودن زمان انتشار توییت، مدل توانایی بیشتری در تحلیل شرایط خاص زمانی بازار به دست می‌آورد و مدلی با نام \lr{Temporal Context-Aware} (\lr{TCA}) توسعه می‌یابد. به عنوان مثال:
\begin{quote}
تاریخ: ۱ ژانویه، ۲۰۲۰ توییت: بیت‌کوین در حال افزایش است!
\end{quote}

از طریق این روش‌ها، فضای ویژگی‌های در دسترس مدل زبانی گسترش یافته و امکان تشخیص روابط پیچیده میان محتوا، شاخص‌های بازار، و دینامیک‌های زمانی فراهم می‌شود که در نتیجه دقت پیش‌بینی در تحلیل مالی افزایش می‌یابد.  این توصیفات به صورت مشخص شده در \ref{fig:tweets} به هر توییت اضافه می‌شود. 

معیارهای ارزیابی برای این آزمایش شامل صحت، دقت، پوشش، امتیاز \lr{F1}،و ماتریس سردرگمی است. این معیارها به طور میانگین از ارزیابی پنج دوره و دو شکن برای هر لایه محاسبه می‌شوند. تمرکز اصلی روی بهبود پیش‌بینی حرکات کوتاه‌مدت بازار با استفاده از برچسب‌گذاری دقیق‌تر نسبت به روش‌های سنتی تحلیل احساسات است.  از ابزار \lr{SHAP} به عنوان تفسیرگر متن برای تحلیل و درک چگونگی تأثیر بخش‌های مختلف توییت بر پیش‌بینی‌های مدل استفاده شد \cite{unified_lundberg_2017}. این تحلیل به منظور بررسی هرگونه نشانه‌ای از بیش‌برازش مدل به بخش‌های خاصی از داده انجام گرفته است. در این بخش، هدف افزودن اطلاعات زمانی به مدل زبان است تا به آن احساس زمان داده شود و قادر به درک این باشد که اطلاعاتی که دریافت می‌کند به زمان خاصی مربوط می‌شود. این کار به امید کاهش بیش‌برازش بیش‌برازش و جلوگیری از گیر افتادن مدل بر روی اطلاعات قبلی انجام می‌شود. برای دستیابی به این هدف، اطلاعات زمانی به هر پیام توییت به صورت غیر توصیفی و با فرمت "سال، ماه، روز" اضافه می‌شود. به این ترتیب، هر توییت با تاریخ و زمان دقیق آن همراه خواهد بود تا مدل بتواند به دقت بیشتری به وابستگی‌های زمانی توجه کند و به درک بهتری از تاثیرات زمانی توییت‌ها دست یابد. این فرآیند به همراه اطلاعات بازار که در بخش قبلی توضیح داده شد، مورد استفاده قرار می‌گیرد. بنابراین، در اینجا به جای استفاده از توضیحات توصیفی برای اطلاعات زمانی، از فرمت تاریخ و زمان استفاده می‌شود که به صورت مستقیم به توییت‌ها افزوده می‌شود. با اضافه کردن این اطلاعات، انتظار می‌رود که مدل بتواند الگوهای زمانی و تغییرات موقتی بازار را بهتر شناسایی کرده و از این طریق، عملکرد بهتری در پیش‌بینی روندهای کوتاه‌مدت بازار داشته باشد. این تغییر به طور خاص در مرحله آموزش مدل با استفاده از روش اعتبارسنجی متقابل گروهی پیاده‌سازی می‌شود، به طوری که هیچ داده‌ای از یک روز خاص در هر دو مجموعه آموزشی و تست وجود ندارد تا از تداخل داده‌ها جلوگیری شود. مدل با استفاده از میانگین رای‌گیری از تمامی فولدها برای هر توییت، پیش‌بینی نهایی را انجام می‌دهد.

\section{تولید سیگنال}
\label{sec:signal-generation}

فرآیند تحلیل احساسات مالی به‌عنوان یکی از ابزارهای پیشرفته در تحلیل بازارهای مالی، هدفی فراتر از طبقه‌بندی ساده توییت‌ها یا متون را دارد. یکی از وظایف کلیدی این فرآیند، تولید سیگنال معاملاتی است که به‌عنوان خروجی نهایی مدل پیشنهادی به کارگزار یا عامل معاملاتی ارائه می‌شود. تولید این سیگنال، در واقع پلی است میان تحلیل احساسات و تصمیم‌گیری در معاملات، که تأثیر مستقیمی بر عملکرد استراتژی‌های معاملاتی دارد. در سیستم‌های مرسوم، مدل‌های مبتنی بر قیمت برای تولید سیگنال‌های معاملاتی استفاده می‌شوند. این مدل‌ها اغلب با مدل‌های تحلیل احساسات ادغام می‌شوند تا ترکیبی از اطلاعات مبتنی بر داده‌های قیمتی و متنی ایجاد شود. به‌عنوان مثال، یک مدل مبتنی بر قیمت ممکن است جهت حرکت بازار را پیش‌بینی کند، در حالی که یک مدل تحلیل احساسات می‌تواند احساسات غالب در بازار را نشان دهد. با این حال، این روش‌ها نیازمند چندین مرحله تحلیل و ترکیب نتایج هستند که می‌تواند منجر به پیچیدگی‌های عملیاتی، تأخیر در تصمیم‌گیری و احتمال افزایش خطا شود. هدف این پژوهش، پیشنهاد یک مدل زبانی تمام است که بتواند تمام این فرآیندها را در یک مرحله ادغام کند. در این رویکرد، مدل زبانی پیشنهادی مستقیماً از داده‌های متنی (مانند توییت‌ها) سیگنال معاملاتی تولید می‌کند و نیازی به مراحل میانی تحلیل قیمت یا ترکیب با سایر مدل‌ها ندارد. این ایده، فرآیند تولید سیگنال را ساده‌تر و یکپارچه‌تر کرده و امکان بهبود سرعت و دقت تصمیم‌گیری معاملاتی را فراهم می‌کند.


\subsection{روش اکثریت}
در روش اکثریت، سیگنال روزانه با تعیین کلاس غالب از میان پیش‌بینی‌های مدل به‌دست می‌آید. فرض کنید تعداد کل توییت‌های پیش‌بینی‌شده در یک روز برابر با \( N \) باشد، و تعداد توییت‌های طبقه‌بندی‌شده به‌عنوان صعودی (\( N_{\text{bullish}} \))، نزولی (\( N_{\text{bearish}} \)) و خنثی (\( N_{\text{neutral}} \)) به ترتیب مشخص شوند. در این روش، سیگنال غالب \( S_{\text{majority}} \) به‌صورت زیر تعریف می‌شود:
\begin{align}
S_{\text{majority}} = \arg\max \left( N_{\text{bullish}}, N_{\text{bearish}}, N_{\text{neutral}} \right)
\end{align}

اطمینان از کلاس غالب، که نشان‌دهنده تسلط آن کلاس در پیش‌بینی‌های روزانه است، به‌صورت زیر محاسبه می‌شود:
\begin{align}
C_{\text{majority}} = \frac{N_{\text{majority}}}{N}
\end{align}
این مقدار نشان‌دهنده‌ی درصدی از توییت‌ها است که به کلاس غالب تعلق دارند. در صورتی که این مقدار بالا باشد، مدل به‌وضوح سیگنال قوی‌تری برای آن روز ارائه داده است. برای اطمینان از مقایسه‌پذیری مقادیر اطمینان در طول زمان، نرمال‌سازی مین-مکس روی مقادیر اطمینان اعمال می‌شود:
\begin{align}
C_{\text{norm}} = \frac{C_{\text{majority}} - \min(C_{\text{majority}})}{\max(C_{\text{majority}}) - \min(C_{\text{majority}})}
\end{align}

\subsection{روش میانگین}
در روش میانگین، سیگنال روزانه از میانگین‌گیری مقادیر کدگذاری‌شده‌ی کلاس‌های پیش‌بینی‌شده محاسبه می‌شود. در این روش، کلاس‌های صعودی، نزولی، و خنثی به ترتیب به‌صورت (نزولی = ۰، خنثی = ۱، و صعودی = ۲) کدگذاری می‌شوند.
فرض کنید \( E_i \) کلاس کدگذاری‌شده‌ی توییت \( i \)-ام باشد. در این صورت، میانگین پیش‌بینی‌های روزانه (\( S_{\text{mean}} \)) به‌صورت زیر محاسبه می‌شود:
\begin{align}
S_{\text{mean}} = \frac{1}{N} \sum_{i=1}^{N} E_i
\end{align}

برای تولید سیگنال، آستانه‌هایی برای تعیین جهت بازار تعریف می‌شوند. این آستانه‌ها شامل \( T_{\text{bearish}} \) برای بازار نزولی و \( T_{\text{bullish}} \) برای بازار صعودی هستند. سیگنال نهایی بر اساس رابطه‌ی زیر تعیین می‌شود:
\begin{align}
S_{\text{mean}} =
\begin{cases} 
\text{نزولی} & \text{اگر } S_{\text{mean}} < T_{\text{bearish}}, \\
\text{صعودی} & \text{اگر } S_{\text{mean}} > T_{\text{bullish}}, \\
\text{خنثی} & \text{در غیر این صورت.}
\end{cases}
\end{align}

اطمینان از سیگنال میانگین، که نشان‌دهنده نزدیکی مقدار میانگین به آستانه‌ها است، به‌صورت زیر محاسبه می‌شود:
\begin{align}
C_{\text{mean}} = S_{\text{mean}} - T(S_{\text{mean}})
\end{align}
برای یکسان‌سازی مقیاس مقادیر اطمینان، این مقادیر نیز با استفاده از رابطه‌ی نرمال‌سازی مین-مکس نرمال می‌شوند. مقادیر اطمینان محاسبه‌شده به عنوان معیاری برای اندازه‌گیری میزان قدرت سیگنال استفاده می‌شوند و در تحلیل‌های بعدی به کار می‌روند.

\subsection{معرفی مدل‌های معیار}
\label{subsec:benchmarks}

به‌منظور انجام تحلیل‌های مقایسه‌ای و ارزیابی عملکرد مدل پیشنهادی، سه مدل معیار طراحی و پیاده‌سازی شدند: 

\begin{enumerate}
    \item \textbf{مدل مبتنی بر قیمت:} یک مدل \lr{LSTM} که از پژوهش \cite{prebit_zou_2023} الهام گرفته شده است. این مدل با استفاده از داده‌های قیمتی و ویژگی‌های مشتق‌شده، پیش‌بینی‌های روزانه‌ای ارائه می‌دهد. ویژگی‌های مورد استفاده در این مدل شامل قیمت‌های باز، بالا، پایین و بسته (\lr{OHLCV})، قیمت اتریوم، قیمت طلا، شاخص قدرت نسبی (\lr{RSI})، نرخ تغییر (\lr{ROC})، میانگین متحرک واگرایی-همگرایی (\lr{MACD}) و میانگین‌های متحرک ساده و نمایی (\lr{MA7، MA21، EMA12، EMA26}) می‌باشند. داده‌های مربوط به سال‌های ۲۰۱۵ تا ۲۰۱۹ برای آموزش این مدل استفاده شده‌اند.
    
    \item \textbf{مدل رمزگذار خودکار:} یک مدل \lr{Autoencoder} که از پژوهش \cite{zha2022time} الهام گرفته شده است. این مدل نیز از همان مجموعه داده و ویژگی‌های قیمتی برای پیش‌بینی‌های روزانه استفاده می‌کند. این مدل به‌طور خاص برای استخراج و کاهش ابعاد ویژگی‌ها طراحی شده و سپس از آن‌ها برای پیش‌بینی استفاده می‌کند.
    
    \item \textbf{مدل ترکیبی\LTRfootnote{Fusion Model}:} این مدل، پیش‌بینی‌های مدل پیشنهادی را با پیش‌بینی‌های مدل‌های مبتنی بر قیمت و پیش‌بینی‌های مبتنی بر احساسات ترکیب می‌کند. پیش‌بینی‌های مبتنی بر احساسات، به صورت نسبت تعداد توییت‌های صعودی به نزولی در طول روز به مدل ترکیبی ارائه می‌شوند. همچنین، پیش‌بینی نهایی مدل‌های قیمتی به‌عنوان یک ویژگی دیگر به این مدل اضافه شده است.
\end{enumerate}

هر یک از این مدل‌های معیار به‌گونه‌ای طراحی شده‌اند که بتوانند نقش و عملکرد بخش‌های مختلف در تولید سیگنال‌های معاملاتی را مورد ارزیابی قرار دهند. مقایسه نتایج این مدل‌ها با مدل پیشنهادی، اطلاعات ارزشمندی درباره مزایا و محدودیت‌های رویکردهای مختلف ارائه می‌دهد.


\section{الگوریتم معاملاتی و پیش‌آزمون}
\label{sec:backtest}
برای فرآیند ارزیابی استراتژی، از ابزار وکتور‌بی‌تی\LTRfootnote{VectorBT} استفاده کردیم. استراتژی پایه‌ای که برای مقایسه استفاده شده است، استراتژی ساده خرید و نگهداری\LTRfootnote{Buy and Hold} است. استراتژی‌های پیشنهادی ما شامل استراتژی برچسب‌گذاری سه‌مانع و استراتژی‌های ورودی-خروجی بر اساس سیگنال‌های مثبت و منفی است. استراتژی سه‌مانع از پارامترهای بهینه‌شده در دامنه‌های خاص برای دوره‌های نوسانات، حدود صعودی و نزولی، و حدود عمودی استفاده می‌کند. این پارامترها به منظور بهبود عملکرد استراتژی تنظیم شده‌اند. از سوی دیگر، استراتژی‌های ورودی-خروجی از سیگنال مخالف به عنوان سیگنال خروج استفاده می‌کنند. سیگنال‌ها برای سه مدل میانگین، اکثریت، و مدل ترکیبی ایجاد شده و برای روش میانگین برچسب روز بر حسب حدود از پیش تعیین‌شده تولید می‌شوند و اندازه معاملات برای هر دو روش میانگین و اکثریت از استاندارد‌سازی بیشترین‌ و کمترین\LTRfootnote{Min-Max Normalization} محاسبه می‌شود تا سطح اطمینان سیگنال را منعکس کند.

معیارهای عملکردی که در طی فرآیند پیش‌آزمون بررسی می‌شوند شامل درصد بازده روزانه، نسبت شارپ، نسبت سورتینو، نسبت سود و افت سرمایه حداکثری هستند. این معیارها ارزیابی جامعی از اثربخشی استراتژی و بازده‌های تنظیم شده بر اساس ریسک ارائه می‌دهند. با مقایسه نتایج با استراتژی پایه خرید و نگهداری، می‌توان عملکرد مدل را در رژیم‌های مختلف بازار ارزیابی کرد. استفاده از پیش‌آزمون به طور مؤثر و کارآمد تضمین می‌کند و امکان تحلیل و مقایسه دقیق استراتژی‌های مختلف را فراهم می‌آورد. این رویکرد بینش‌های ارزشمندی در مورد نقاط قوت و ضعف هر استراتژی فراهم می‌آورد و به بهبودها و کاربردهای آینده در پیش‌بینی بازارهای مالی راهنمایی می‌کند.

\subsection{پیش‌آزمون استراتژی‌شده}

برای ارزیابی عملکرد مدل‌های زبانی در محیط معامله، از طریق سیگنال‌های به دست آمده از مدل تنظیم‌شده بر برچسب‌گذاری پیشنهاد‌ی، به پیش‌آزمون این استراتژی‌ها میپردازیم. فرضیه اصلی این است که مدل‌های زبانی که برای پیش‌بینی کوتاه‌مدت بازار آموزش دیده‌اند، در کنار روش‌های معاملاتی تعریف‌شده، عملکرد قابل‌توجهی خواهند داشت. داده‌های مورد استفاده شامل سه دوره زمانی یک‌ساله که هر یک نماینده یکی از رژیم‌های صعودی، نزولی، و خنثی بازار هستند. این تقسیم‌بندی به منظور ارزیابی عملکرد مدل‌ها در شرایط مختلف بازار انجام شده است. برای جلوگیری از انحراف در نتایج، نمونه‌برداری مجدد انجام نشده و داده‌ها به صورت یکپارچه در هر دوره استفاده شده‌اند. سیگنال‌های معاملاتی بر اساس این داده‌ها ایجاد و در استراتژی‌های مختلف به کار گرفته شده‌اند. معیارهای ارزیابی شامل درصد بازده روزانه، نسبت شارپ، نسبت سورتینو، و کاهش حداکثری سرمایه است. نتایج حاصل از پیش‌آزمون با استراتژی پایه خرید و نگهداری و استراتژی‌های ورود و خروج مقایسه شده‌اند تا اثربخشی سیگنال‌های مدل در رژیم‌های مختلف بازار ارزیابی شود. این تحلیل به شناسایی نقاط قوت و ضعف استراتژی‌های پیشنهادی و ارائه بینش‌هایی برای بهبود و توسعه روش‌های پیش‌بینی در آینده کمک می‌کند.
\subsection{استراتژی‌های معاملاتی}
اینجا به سه استراتژی استفاده شده در محیط پیش‌آزمون میپردازیم:
\begin{itemize}
    \item استراتژی خرید و نگه‌داری:
     استراتژی خرید و نگه‌داری شامل خرید در شروع بازه و فروش در پایان بازه می‌شود. با توجه به وابستگی شدید این نوع معامله در بازار، برای بازه صعودی و خنثی در بازار از خرید و نگه‌داری استفاده می‌کنیم و برای بازه نزولی از فروش و نگه‌داری استفاده می‌کنیم. انتظار داریم که با توجه به درصد بالای بازدهی بازار در بازه صعودی، عملکرد خوبی از این استراتژی به دست بیاید و به عنوان معیار مقایسه از آن استفاده می‌کنیم. در این استراتژی، تنها یک موقعیت باز در هر زمان وجود دارد و معامله‌گر تا پایان بازه زمانی تعیین شده، موقعیت خود را نگه می‌دارد.
    
    \item استراتژی ورود و خروج:
    این استراتژی در قالب دو حالت ورود و خروج صعودی و ورود و خروج نزولی استفاده شده است:
    \begin{itemize}
        \item ورود و خروج صعودی:
        در این نوع از استراتژی، با سیگنال صعودی وارد معامله می‌شویم و با سیگنال عکس (نزولی) خارج می‌شویم. تنها یک موقعیت باز در هر زمان وجود دارد و معامله‌گر با دریافت سیگنال صعودی وارد بازار می‌شود و با دریافت سیگنال نزولی از بازار خارج می‌شود.
        
        \item ورود و خروج نزولی:
        در این نوع از استراتژی، با سیگنال نزولی وارد معامله فروش\LTRfootnote{Short} می‌شویم و با سیگنال صعودی خارج می‌شویم. تنها یک موقعیت باز در هر زمان وجود دارد و معامله‌گر با دریافت سیگنال نزولی وارد بازار می‌شود و با دریافت سیگنال صعودی از بازار خارج می‌شود.
    \end{itemize}
    
    \item استراتژی سه‌مانع:
    این استراتژی با استفاده از سیگنال‌های ورود وارد بازار شده و با اصابت به حدود پویای تعیین شده توسط نوسان معامله را به پایان می‌رساند. در این روش، از برچسب‌های سه‌مانع برای تعیین نقاط ورود و خروج استفاده می‌شود تا به دقت بیشتری در پیش‌بینی حرکات بازار دست یابیم. تنها یک موقعیت باز در هر زمان وجود دارد و معامله‌گر با دریافت سیگنال ورود وارد بازار می‌شود و با رسیدن به حدود تعیین شده توسط نوسان از بازار خارج می‌شود.
\end{itemize}

استراتژی‌های معاملاتی تعریف‌شده در این بخش با هدف ارزیابی اعتبار و کارایی روش پیشنهادی در شرایط مختلف بازار، شامل رژیم‌های صعودی، نزولی، و خنثی طراحی شده‌اند. با بررسی عملکرد این استراتژی‌ها در دوره‌های زمانی مشخص، امکان تحلیل تطبیق‌پذیری مدل‌های زبانی و سیگنال‌های تولیدشده آن‌ها با وضعیت‌های متغیر بازار فراهم می‌شود. این رویکرد به درک عمیق‌تر از توانایی مدل‌های پیشنهادی در پیش‌بینی حرکات بازار و مدیریت موقعیت‌های معاملاتی کمک می‌کند.
